\chapter{Conclusiones}
\label{cap:conclusiones}

% ===================================================================
% Pautas de redacción (pautas.md), bloque «la conclusión»:
%   1. Resumir las implicaciones de los resultados, tal como se enunciaron en
%      las discusiones de los capítulos 4, 5 y 6. Sin cifras nuevas.
%   2. Describir cómo se alcanzaron los objetivos.
%   3. Dejar explícito, sin que el lector lo deduzca: si se cumplió cada
%      resultado esperado, cada objetivo específico, el objetivo general, y si
%      la solución resuelve el problema central.
%
% El capítulo es corto a propósito: las cinco preguntas de la pauta de
% resultados (significado, literatura, generalización, limitaciones) ya están
% respondidas en la Discusión de cada capítulo de objetivo, y repetirlas aquí
% sería redundancia.
%
% BLOQUEADO POR EL CAPÍTULO 6. Las secciones de OE3 y del objetivo general no
% pueden redactarse hasta que RE3.2 entregue los minutos por hora de audio en
% condición manual y asistida: el OG se compromete con esa magnitud y no con
% ninguna otra. Lo que falta está señalado abajo en cada sección.
% ===================================================================

\section{Implicaciones de los resultados}

El trabajo deja dos cosas que no existían al empezar. La primera es un conjunto
acústico utilizable: \nKept{} anotaciones curadas bajo un vocabulario cerrado, de
las que \nExperiment{} forman un experimento de \nClasses{} clases particionado
por grabación, con su ficha y sus limitaciones declaradas. La segunda es la
evidencia de que la detección de cajas tiempo--frecuencia sobre espectrograma es
viable en este dominio, y de que el cuello de botella no está donde se esperaba.

Esa es la implicación principal del Capítulo~\ref{cap:deteccion}: las cuatro
configuraciones comparadas no se separan en la métrica primaria, con intervalos
de confianza solapados y un rango de \rRecallSpread{} en \textit{recall}, mientras
que la clasificación de especie y tipo de llamada resulta casi resuelta sobre los
pares emparejados. El problema difícil no era distinguir \nClasses{} clases con
dos órdenes de magnitud de desbalance, sino encontrar y encuadrar el evento. Para
el equipo de investigación la consecuencia es concreta: el analista no tendrá que
corregir etiquetas, tendrá que corregir cajas.

La segunda implicación es metodológica. Al declarar la regla de decisión antes de
entrenar, el experimento pudo cerrarse con un resultado que no favorece a la
arquitectura propuesta sin que eso lo invalide, y al separar la medición en
detección, encuadre y clasificación quedó visible qué parte del problema está
resuelta y cuál no. Una sola cifra agregada habría escondido ambas cosas.

\section{Cumplimiento de los resultados esperados}

La Tabla~\ref{tab:cumplimiento} contrasta cada resultado esperado del
Capítulo~\ref{cap:generalidades} con la sección que lo entrega y lo verifica
contra su indicador.

\begin{table}[htbp]
    \centering
    \caption{Cumplimiento de los resultados esperados}
    \label{tab:cumplimiento}
    \small
    \begin{tabular}{@{}lL{7.2cm}L{3.4cm}l@{}}
        \toprule
        \textbf{RE} & \textbf{Resultado} & \textbf{Verificado en} & \textbf{Estado} \\
        \midrule
        RE1.1 & Protocolo de curación documentado        & Sección~\ref{sec:verificacion-re11} & Cumplido \\
        RE1.2 & Conjunto curado y particionado           & Sección~\ref{sec:verificacion-re12} & Cumplido \\
        RE1.3 & Ficha del conjunto derivado              & Sección~\ref{sec:verificacion-re13} & Cumplido \\
        \addlinespace
        RE2.1 & Diseño, arquitectura y líneas base       & Sección~\ref{sec:verificacion-re21} & Cumplido \\
        RE2.2 & \textit{Pipeline} reproducible           & Sección~\ref{sec:verificacion-re22} & Cumplido \\
        RE2.3 & Informe comparativo desagregado          & Sección~\ref{sec:verificacion-re23} & Cumplido \\
        RE2.4 & Modelo final y salida a tabla de Raven   & Sección~\ref{sec:demo-oe2} & Cumplido \\
        RE2.5 & Repositorio con código y modelo          & Sección~\ref{sec:re25} & Cumplido \\
        \bottomrule
    \end{tabular}
\end{table}

% PENDIENTE (capítulo 6): añadir a la tabla las filas de RE3.1, RE3.2 y RE3.3 con
% su sección de verificación y su veredicto, y ampliar el párrafo siguiente.

Los ocho resultados de OE1 y OE2 se entregan y se verifican contra su indicador.
Uno lo hace con una desviación declarada en la sección correspondiente, no
silenciada: la dimensión del decodificador y el número de consultas de la
arquitectura propuesta difieren de los valores anunciados en RE2.1, y la
sección~\ref{sec:verificacion-re21} razona por qué.

\section{Cumplimiento de los objetivos específicos}

\rotulo{OE1.} Cumplido. El Capítulo~\ref{cap:curacion} entrega el protocolo, el
conjunto particionado por grabación y la ficha, y mide cada uno contra su
indicador: el protocolo conserva el \pctKept{} de las filas crudas, la partición
no comparte ninguna grabación entre subconjuntos y la ficha cubre los apartados
comprometidos, incluidas las limitaciones que el propio material impone.

\rotulo{OE2.} Cumplido. El Capítulo~\ref{cap:deteccion} implementa el detector,
lo compara con tres líneas base bajo un protocolo único declarado de antemano y
selecciona \rWinner{} por la métrica secundaria, con \rRecallYol{} de
\textit{recall} y \rMapYol{} de mAP en la partición de prueba, medida una sola
vez. El modelo final se carga y produce una tabla que Raven abre sobre la
grabación, y el repositorio contiene el código y el \textit{checkpoint}.

% PENDIENTE (capítulo 6): el párrafo de OE3, con el mismo formato, una vez que
% RE3.1 a RE3.3 estén entregados y verificados.

\section{Cumplimiento del objetivo general}

% PENDIENTE (capítulo 6). El objetivo general se compromete con la reducción del
% tiempo de revisión experta, y esa magnitud sale de RE3.2. Cuando exista, el
% veredicto se enuncia aquí sin rodeos: cuántos minutos por hora de audio cuesta
% la revisión manual, cuántos la asistida, y en qué condiciones vale la
% comparación. Hasta entonces no hay respuesta que dar, y adelantarla sería
% exactamente lo que el Capítulo 5 evitó al declarar la regla de decisión antes
% de medir.

\section{La propuesta frente al problema central}

El problema central del Capítulo~\ref{cap:generalidades} tiene dos predicados: se
anota mucho menos de lo que se graba, y lo que se anota no obedece a un criterio
único. La alternativa de solución ataca los dos, y lo hace por vías distintas.

Contra el déficit de consistencia actúa OE1. El vocabulario cerrado, el mapa de
sinónimos y la regla que hace de la carpeta la fuente de verdad de la especie
convierten un material escrito a mano, con decenas de grafías para la misma
clase, en \nClasses{} clases con criterio único. No corrige el juicio del
anotador, que es irrecuperable sin volver al audio, pero elimina la parte del
desacuerdo que era de escritura y no de criterio, y deja marcado y conservado lo
que no puede decidir.

Contra el déficit de caudal actúa OE2, y aquí conviene ser preciso sobre lo que el
sistema es y lo que no. No sustituye al analista ni pretende hacerlo: la cabeza
emite cajas con la forma exacta en que el equipo ya anota, y el analista las
recibe en su propia herramienta para aceptarlas, corregirlas o borrarlas. Es
pre-anotación asistida, y por eso el compromiso elegido privilegia la cobertura
sobre la precisión: perder una vocalización cuesta más que descartar una
propuesta falsa, porque el humano revisa la salida de todos modos.

% PENDIENTE (capítulo 6). Lo anterior establece que la propuesta se enfoca en el
% problema central, que es lo que pide la pauta. Cuánto lo reduce es otra
% pregunta, y la responde RE3.2: hasta tenerla, este capítulo no afirma que el
% caudal de anotación haya aumentado, sólo que el sistema está construido para
% aumentarlo y que entrega la salida en la forma que ese aumento requiere.

\section{Trabajo futuro}

Lo que queda fuera no es una lista de deseos: son las limitaciones declaradas en
las discusiones de los capítulos de objetivo, ordenadas por cuánto comprometen lo
que aquí se afirma.

\begin{itemize}
    \item \textbf{Medir el acuerdo entre anotadores.} Es la carencia más seria.
          Sin ella no hay techo de evaluación, y una parte del error que este
          trabajo atribuye al modelo podría ser desacuerdo humano.
    \item \textbf{Separar el extractor de la cabeza.} Las tres ablaciones que lo
          harían están implementadas y sin correr, de modo que la comparación
          entre configuraciones no puede atribuir la diferencia a un componente.
    \item \textbf{Reconstruir el evento largo.} Las clases de duración mayor que
          la ventana sobreviven fragmentadas, sin ningún mecanismo que las
          reagrupe antes de puntuarlas.
    \item \textbf{Cubrir el nivel de frase.} Las clases de frase quedan fuera por
          anidamiento jerárquico, y con ellas el nivel superior de la jerarquía
          vocal.
    \item \textbf{Cerrar el vocabulario con el repertorio de referencia.} Los
          pares fuera de vocabulario mezclan error tipográfico y subdivisión
          legítima, y no hay forma de separarlos sin una referencia del
          repertorio vocal de cada especie.
\end{itemize}
